% BHU PYQ -- AIHC & Archaeology: Ancient Indian Social Life and Institutions (2025-26 NEP)
\documentclass[11pt,a4paper]{article}
\usepackage[a4paper,top=2.5cm,bottom=2.5cm,left=2.8cm,right=2.8cm,headheight=14pt]{geometry}
\usepackage{amsmath,amssymb}
\usepackage{fontspec}
\setmainfont{Kohinoor Devanagari}
\usepackage{microtype,fancyhdr,enumitem,hyperref,lastpage,parskip}

\hypersetup{colorlinks=true,urlcolor=black,linkcolor=black,pdftitle={AIHMJ32: Ancient Indian Social Life and Institutions -- BHU NEP 2025-26}}

\pagestyle{fancy}\fancyhf{}
\renewcommand{\headrulewidth}{0.4pt}\renewcommand{\footrulewidth}{0.4pt}
\fancyhead[L]{\small \textit{Banaras Hindu University}}
\fancyhead[R]{\small \textit{AIHMJ32: Social Life and Institutions (2025-26)}}
\fancyfoot[C]{\small Page \thepage\ of \pageref{LastPage}}

\newlist{parts}{enumerate}{2}
\setlist[parts,1]{label=(\alph*),leftmargin=*,itemsep=4pt,topsep=3pt}
\setlist[parts,2]{label=(\roman*),leftmargin=*,itemsep=2pt,topsep=2pt}

\newcommand{\pts}[1]{\hfill \textit{[#1]}}

\begin{document}

\begin{center}
  {\large\bfseries BANARAS HINDU UNIVERSITY}\\[2pt]
  {\normalsize B. A. (Hons.) Semester III Examination, 2025-26 (NEP)}\\[6pt]
  \rule{0.8\linewidth}{0.5pt}\\[6pt]
  {\large\bfseries ANCIENT INDIAN HISTORY, CULTURE \& ARCHAEOLOGY}\\[4pt]
  {\large\bfseries Paper Code: AIHMJ32 : Ancient Indian Social Life and Institutions}\\[4pt]
  {\normalsize Time: 2:30 Hours \hfill Full Marks: 70}\\[2pt]
  {\small REF NO. 2025-26/D-III/069}
\end{center}
\rule{\linewidth}{0.4pt}
\smallskip

\noindent \textit{Instructions: Attempt all questions of Sections A, B and C which carry 10, 5 and 2 marks each respectively. Write your Roll No.\ at the top immediately on receipt of this question paper.}

\smallskip
\noindent \textit{खण्ड अ, ब तथा स के सभी प्रश्नों के उत्तर दीजिये जिनके अंक क्रमशः 10, 5 तथा 2 हैं।}

\bigskip

\section*{SECTION -- A / खण्ड -- अ}
\noindent \textit{Note: Long Answer Type Questions. Answer approximately in 500--700 words.}\pts{3 $\times$ 10 = 30}

\noindent \textit{दीर्घ उत्तरीय प्रश्न। लगभग 500--700 शब्दों में उत्तर दीजिये।}

\begin{enumerate}[label=\textbf{\arabic*.},leftmargin=*,itemsep=12pt]

  \item Describe the various theories of the origin of the Varna system.\pts{10}
  
  वर्ण व्यवस्था की उत्पत्ति के विभिन्न सिद्धान्तों का वर्णन कीजिये।

  \begin{center}\textbf{OR / अथवा}\end{center}

  Throw light on the condition of women in ancient India.\pts{10}
  
  प्राचीन भारत में नारियों की स्थिति पर प्रकाश डालिये।

  \item Write a note on the Ashrama system of ancient India.\pts{10}
  
  प्राचीन भारत की आश्रम व्यवस्था पर एक टिप्पणी लिखिये।

  \begin{center}\textbf{OR / अथवा}\end{center}

  Explain the meaning of Caste and throw light on its complexities.\pts{10}
  
  जाति के अर्थ को स्पष्ट कीजिये तथा इसकी जटिलताओं पर प्रकाश डालिये।

  \item Highlight the educational centres of Takshashila and Nalanda, bringing out the importance of ancient Indian education.\pts{10}
  
  प्राचीन भारतीय शिक्षा के महत्त्व को रेखांकित करते हुए तक्षशिला एवं नालंदा शिक्षा केंद्रों पर प्रकाश डालिये।

  \begin{center}\textbf{OR / अथवा}\end{center}

  Throw light on the Sixteen Samskaras.\pts{10}
  
  षोडश (सोलह) संस्कारों पर प्रकाश डालिये।

\end{enumerate}

\bigskip

\section*{SECTION -- B / खण्ड -- ब}
\noindent \textit{Note: Short Answer Type Questions. Answer approximately in 250 words.}\pts{4 $\times$ 5 = 20}

\noindent \textit{लघु उत्तरीय प्रश्न। लगभग 250 शब्दों में उत्तर दीजिये।}

\begin{enumerate}[label=\textbf{\arabic*.},leftmargin=*,start=4,itemsep=10pt]

  \item Throw light on the status of Shudras in ancient Indian society.\pts{5}
  
  प्राचीन भारतीय समाज में शूद्रों की स्थिति पर प्रकाश डालिये।

  \begin{center}\textbf{OR / अथवा}\end{center}

  Throw light on the concept of Rina (indebtedness) prevalent in ancient India.\pts{5}
  
  प्राचीन भारत में प्रचलित 'ऋण' की अवधारणा पर प्रकाश डालिये।

  \item What is marriage? Explain its different types.\pts{5}
  
  विवाह क्या है? इसके विभिन्न प्रकारों की व्याख्या कीजिये।

  \begin{center}\textbf{OR / अथवा}\end{center}

  Explain the concept of family and highlight the importance of parents.\pts{5}
  
  परिवार की अवधारणा को स्पष्ट कीजिये तथा माता-पिता के महत्त्व पर प्रकाश डालिये।

  \item Throw light on the types of sons (Putra) in ancient Indian society.\pts{5}
  
  प्राचीन भारतीय समाज में पुत्रों के प्रकारों पर प्रकाश डालिये।

  \begin{center}\textbf{OR / अथवा}\end{center}

  What do you understand by Purushartha? Describe any two Purusharthas.\pts{5}
  
  पुरुषार्थ से आप क्या समझते हैं? किन्हीं दो पुरुषार्थों का वर्णन कीजिये।

  \item Write an essay on the status of widows in ancient Indian society.\pts{5}
  
  प्राचीन भारतीय समाज में विधवाओं की स्थिति पर एक निबन्ध लिखिये।

  \begin{center}\textbf{OR / अथवा}\end{center}

  Bring out the significance of the Upanayana ceremony in ancient India.\pts{5}
  
  प्राचीन भारत में उपनयन संस्कार के महत्त्व को रेखांकित कीजिये।

\end{enumerate}

\bigskip

\section*{SECTION -- C / खण्ड -- स}
\noindent \textit{Note: Very Short Answer Type Questions. Answer in the briefest possible way.}\pts{10 $\times$ 2 = 20}

\noindent \textit{अति संक्षिप्त उत्तर दीजिये।}

\begin{enumerate}[label=\textbf{\arabic*.},leftmargin=*,start=8,itemsep=8pt]
  \item 
  \begin{parts}
    \item Who was called `Punarbhu'?\pts{2}
    
    `पुनर्भू' किसे कहा जाता था? (Remarried widow / पुनर्विवाहिता नारी)

    \item Which Ashrama has Manu called the best Ashrama?\pts{2}
    
    मनु ने किस आश्रम को श्रेष्ठतम आश्रम कहा है? (Grihastha Ashrama / गृहस्थाश्रम)

    \item What do you understand by Upanishad?\pts{2}
    
    उपनिषद् से आप क्या समझते हैं?

    \item Mention any two types of slaves (Dasa) in ancient India.\pts{2}
    
    प्राचीन भारत में दासों के किन्हीं दो प्रकारों का उल्लेख कीजिये। (Dhvajahrita, Garbhadasa)

    \item Who is the author of the book ``Hindu Samskara''?\pts{2}
    
    ``हिन्दू संस्कार'' पुस्तक के लेखक कौन हैं? (Dr. Rajbali Pandey)

    \item What do you understand by `Purusha Sukta'?\pts{2}
    
    `पुरुष सूक्त' से आप क्या समझते हैं? (Hymn in 10th Mandala of Rigveda describing 4 Varnas)

    \item What do you mean by `Dwij'?\pts{2}
    
    `द्विज' से आप क्या समझते हैं? (Twice-born: Brahmana, Kshatriya, Vaishya after Upanayana)

    \item What are the four ancient marriages of the highest order (Shasta/Prashasta)?\pts{2}
    
    धर्मसम्मत/प्रशस्त चार प्राचीन विवाह कौन-से हैं? (Brahma, Daiva, Arsha, Prajapatya)

    \item Who was called `Brahmavadini'?\pts{2}
    
    `ब्रह्मवादिनी' किसे कहा जाता था? (Women devoted to lifelong Vedic study)

    \item What do you understand by `Deva Rina'?\pts{2}
    
    `देव ऋण' से आप क्या समझते हैं? (Debt to Gods, cleared through Yajna/sacrifices)
  \end{parts}
\end{enumerate}

\end{document}
